\textbf{1. Kiến thức cần nhớ}
\begin{itemize}
\item \textbf{Hình thang} là tứ giác có hai cạnh đối song song ( \parallel CD$). Hai cạnh song song gọi là \textit{hai đáy}, hai cạnh còn lại gọi là \textit{hai cạnh bên}.
\item \textbf{Hình thang cân} là hình thang có hai góc kề một đáy bằng nhau.
\item \textbf{Tính chất hình thang cân}: Trong hình thang cân:
\begin{itemize}
\item Hai cạnh bên bằng nhau: $\boxed{AD = BC}$.
\item Hai đường chéo bằng nhau: $\boxed{AC = BD}$.
\end{itemize}
\item \textbf{Dấu hiệu nhận biết hình thang cân}:
\begin{enumerate}
\item Hình thang có hai góc kề một đáy bằng nhau là hình thang cân.
\item Hình thang có hai đường chéo bằng nhau là hình thang cân.
\end{enumerate}
\item \textbf{Lưu ý}: Tứ giác có hai cạnh bên bằng nhau chưa chắc là hình thang cân (có thể là hình bình hành).
\end{itemize}

\textbf{2. Phương pháp giải các dạng}

%% DANG 8H3?11-1
\textbf{Dạng 1. Tính chất hình thang, hình thang cân}
\begin{enumerate}
\item Vận dụng tính chất hai đáy song song: hai góc kề một cạnh bên bù nhau (tổng bằng ^\circ$, do ở vị trí trong cùng phía).
\item Sử dụng tính chất hai góc kề một đáy bằng nhau, hai cạnh bên bằng nhau, hai đường chéo bằng nhau để suy ra các đoạn thẳng, góc tương ứng bằng nhau.
\end{enumerate}

%% DANG 8H3?11-2
\textbf{Dạng 2. Dấu hiệu nhận biết hình thang cân}
\begin{enumerate}
\item \textbf{Bước 1}: Chứng minh tứ giác là hình thang (chỉ ra hai cạnh đối song song nhờ cặp góc so le trong bằng nhau hoặc đồng vị bằng nhau).
\item \textbf{Bước 2}: Chứng minh hình thang đó có một trong hai dấu hiệu:
\begin{itemize}
\item Hai góc kề một đáy bằng nhau ($\widehat{C} = \widehat{D}$ hoặc $\widehat{A} = \widehat{B}$).
\item Hai đường chéo bằng nhau ( = BD$).
\end{itemize}
\end{enumerate}

%% DANG 8H3?11-3
\textbf{Dạng 3. Tính độ dài, số đo góc trong hình thang cân}
\begin{enumerate}
\item Kẻ thêm hai đường cao từ hai đỉnh của đáy nhỏ xuống đáy lớn ( \perp CD, BK \perp CD$).
\item Khi đó $ là hình chữ nhật nên  = AB$. Hai tam giác vuông $\triangle AHD$ và $\triangle BKC$ bằng nhau (cạnh huyền - góc nhọn) suy ra:
80669\boxed{DH = CK = \dfrac{CD - AB}{2}}80669
\item Áp dụng định lí Pythagore trong tam giác vuông $ để tính chiều cao $ hoặc đường chéo $.
\end{enumerate}
